\ifdefined\gkver\else\def\gkver{student}\fi
\documentclass[\gkver]{gaokaozhenti}
\begin{document}

\chapter{2026年陕晋青宁}

\begin{enumerate}
\item
%% number: 1
%% paperName: 2026年陕晋青宁高考第 1 题 4 分
%% typeId: 1
%% type: 单选题
%% chapter: 直线运动
%% point: 路程与位移
%% method: 无
%% score: 4
%% degree: 950
%% duplicateId: 0
%% body:
现代科技为物资运输提供了多种方式。某物流企业分别利用无人机和新能源货车从山上装货点 P 运输货物至山下仓库 Q，两种方式中一定相同的是 \xzanswer{A}

\fourchoices[answer=A]
{位移}
{路程}
{时间}
{平均速度}

%% answer: A
%% memo:
\memoanswer{
A．位移是从初位置指向末位置的有向线段，仅由初、末位置决定，两种运输方式的初位置都是 P、末位置都是 Q，故位移一定相同，A 正确；

B．路程是物体运动轨迹的长度，无人机和货车的行驶路径不同，运动轨迹长度不一定相等，故路程不一定相同，B 错误；

C．运动时间由运动路程和运行速率共同决定，两者的路程、运行速率均无确定关系，故时间不一定相同，C 错误；

D．平均速度的定义为 $\overline{v}=\frac{x}{t}$，两者位移 $x$ 相同，但运动时间 $t$ 不一定相同，故平均速度不一定相同，D 错误。

故选 A。
}

\item
%% number: 2
%% paperName: 2026年陕晋青宁高考第 2 题 4 分
%% typeId: 1
%% type: 单选题
%% chapter: 力物体的平衡
%% point: 多力平衡
%% method: 无
%% score: 4
%% degree: 850
%% duplicateId: 0
%% body:
我国科研人员利用仿生机器鱼研究湖泊生态，当仿生机器鱼在湖中匀速直线下潜时，它受到水的合力方向是 \xzanswer{D}

\fourchoices[answer=D]
{斜向下}
{竖直向下}
{斜向上}
{竖直向上}

%% answer: D
%% memo:
\memoanswer{
仿生机器鱼做匀速直线运动，处于平衡状态，所受合外力为 0。机器鱼受两个外力：竖直向下的重力 $G$，以及水对其施加的所有作用力的合力 $F$（包含浮力、水的阻力等）。根据平衡条件有 $G+F=0$，可得 $F=-G$，

即水的合力与重力大小相等、方向相反，为竖直向上。

故选 D。
}

\item
%% number: 3
%% paperName: 2026年陕晋青宁高考第 3 题 4 分
%% typeId: 1
%% type: 单选题
%% chapter: 交流电
%% point: 交流电
%% method: 无
%% score: 4
%% degree: 850
%% duplicateId: 0
%% body:
下列处于匀强磁场的导体回路中能产生交变电流的是 \xzanswer{B}

\fourpicture[numstyle=chinese, labelprefix={}, h=3.0cm]
{figs/03a.pdf}
{figs/03b.pdf}
{figs/03c.pdf}
{figs/03d.pdf}

\fourchoices[answer=B]
{甲图中绕轴匀速转动的导体环}
{乙图中绕轴匀速转动的导体框}
{丙图中在固定导轨上匀速运动的导体棒所在回路}
{丁图中匀速运动的导体框}

%% answer: B
%% memo:
\memoanswer{
A．甲：导体环绕轴匀速转动，磁通量不改变，不产生交变电流，故 A 错误；

B．乙：导体框绕轴匀速转动，$\Phi=BS\cos\omega t$，$\frac{\Delta\Phi}{\Delta t}$ 周期性变化，产生交变电流，故 B 正确；

C．丙：导体棒匀速运动，感应电动势 $E=BLv$，电流恒定，故 C 错误；

D．丁：导体框匀速运动，$\Phi$ 不变，无感应电流，故 D 错误。

故选 B。
}

\item
%% number: 4
%% paperName: 2026年陕晋青宁高考第 4 题 4 分
%% typeId: 1
%% type: 单选题
%% chapter: 运动和力
%% point: 动量和能量
%% method: 无
%% score: 4
%% degree: 650
%% duplicateId: 0
%% body:
正电子 $e^{+}$ 与电子 $e^{-}$ 湮灭产生的光子是探测中微子的重要信号。一个静止的正电子与一个静止的电子湮灭转化为两个 $\gamma$ 光子，该过程动量和能量都守恒。已知正电子与电子的质量均为 $m$，光速为 $c$，普朗克常量为 $h$，则这两个 $\gamma$ 光子 \xzanswer{B}

\fourchoices[answer=B]
{总能量为 $mc^{2}$}
{频率均为 $\frac{mc^{2}}{h}$}
{总动量为 $2mc$}
{波长均为 $\frac{mc}{h}$}

%% answer: B
%% memo:
\memoanswer{
静止的正电子和电子总静能为 $2mc^{2}$，总动量为 0，湮灭过程满足能量守恒和动量守恒，因此两个光子总能量为 $2mc^{2}$，总动量为 0，两光子反向运动，每个光子的能量、动量大小均相等，单个光子能量为 $mc^{2}$，动量大小 $p=\frac{E}{c}=mc$，

A．两个光子总能量等于正负电子总静能 $2mc^{2}$，并非 $mc^{2}$，故 A 错误；

B．单个光子能量满足 $E=h\nu=mc^{2}$，解得频率 $\nu=\frac{mc^{2}}{h}$，故 B 正确；

C．初始系统总动量为 0，由动量守恒可知两个光子总动量为 0，并非 $2mc$，故 C 错误；

D．光子动量满足 $p=\frac{h}{\lambda}=mc$，

解得波长 $\lambda=\frac{h}{mc}$，并非 $\frac{mc}{h}$，故 D 错误。

故选 B。
}

\item
%% number: 5
%% paperName: 2026年陕晋青宁高考第 5 题 4 分
%% typeId: 1
%% type: 单选题
%% chapter: 曲线运动
%% point: 人造地球卫星
%% method: 无
%% score: 4
%% degree: 550
%% duplicateId: 0
%% body:
我国自主建设运行的北斗全球卫星导航系统空间段由多颗卫星组成。若轨道半径不同的两颗卫星绕地球做匀速圆周运动，角速度大小分别为 $\omega_{1}$、$\omega_{2}$，线速度大小分别为 $v_{1}$、$v_{2}$，则 \xzanswer{D}

\fourchoices[answer=D]
{$v_{1}/\omega_{1}=v_{2}/\omega_{2}$}
{$v_{1}/\omega_{1}^{2}=v_{2}/\omega_{2}^{2}$}
{$v_{1}^{2}/\omega_{1}=v_{2}^{2}/\omega_{2}$}
{$v_{1}^{3}/\omega_{1}=v_{2}^{3}/\omega_{2}$}

%% answer: D
%% memo:
\memoanswer{
卫星绕地球做匀速圆周运动，万有引力提供向心力 $G\frac{Mm}{r^{2}}=m\omega^{2}r=m\frac{v^{2}}{r}$，

其中 $G$ 为引力常量，$M$ 为地球质量，$r$ 为轨道半径，$\omega$ 为角速度，$v$ 为线速度。

推导得 $v=\sqrt{\frac{GM}{r}}$，$\omega=\sqrt{\frac{GM}{r^{3}}}$，

A．由 $v=\omega r$，由于轨道半径不同，则 $v_{1}/\omega_{1}\neq v_{2}/\omega_{2}$，故 A 错误；

B．由以上分析可知，$\frac{v}{\omega^{2}}\propto r^{\frac{5}{2}}$，与 $r$ 有关，$v_{1}/\omega_{1}^{2}\neq v_{2}/\omega_{2}^{2}$，故 B 错误；

C．由以上分析可知，$\frac{v^{2}}{\omega}\propto\sqrt{r}$，与 $r$ 有关，$v_{1}^{2}/\omega_{1}\neq v_{2}^{2}/\omega_{2}$，故 C 错误；

D．由以上分析可知，$\frac{v^{3}}{\omega}=\frac{\left(\frac{GM}{r}\right)^{\frac{3}{2}}}{\sqrt{\frac{GM}{r^{3}}}}=GM$，与 $r$ 无关，故 D 正确。

故选 D。
}

\item
%% number: 6
%% paperName: 2026年陕晋青宁高考第 6 题 4 分
%% typeId: 1
%% type: 单选题
%% chapter: 电磁感应
%% point: 电磁感应中的变加速
%% method: 无
%% score: 4
%% degree: 450
%% duplicateId: 0
%% body:
电磁泵减速部分的简化工作原理如图。在垂直纸面向外、磁感应强度大小为 $B$ 的匀强磁场中，竖直放置的 U 形光滑金属导轨左右间距为 $L$，水平放置的导体棒 ab 质量为 $m$，与竖直放置且下端固定的绝缘轻弹簧相连，ab 与导轨接触良好且其接入回路的电阻阻值为 $R$。将 ab 自弹簧原长位置由静止释放，ab 第一次到达最低点时下降高度为 $h$，此过程中其最大速度为 $v$。导轨电阻忽略不计，弹簧始终在弹性限度内。在 ab 从静止释放到第一次到达最低点的过程中 \xzanswer{C}

\onepicture[h=4.2cm]{\includesvg{figs/06}}

\fourchoices[answer=C]
{流过导体棒的感应电流方向由 a 向 b}
{导体棒所受安培力的冲量为零}
{通过导体棒的电荷量大小为 $\frac{BLh}{R}$}
{导体棒的加速度最大值为 $\frac{B^{2}L^{2}v}{mR}$}

%% answer: C
%% memo:
\memoanswer{
A．由右手定则可知，流过导体棒的感应电流方向由 b 向 a，故 A 错误；

B．导体棒下降过程中，安培力大小改变，方向不变，由冲量公式 $I=Ft$ 可知，导体棒所受安培力的冲量不为零，故 B 错误；

C．导体棒下降过程中，感应电动势为 $\overline{E}=\frac{\Delta\Phi}{\Delta t}=\frac{BLh}{\Delta t}$，

感应电流为 $\overline{I}=\frac{\overline{E}}{R}=\frac{BLh}{R\cdot\Delta t}$，

通过导体棒的电荷量大小为 $q=\overline{I}\cdot\Delta t=\frac{BLh}{R}$，故 C 正确；

D．根据题意可知，导体棒下降过程中，先做加速度减小的加速运动，再做加速度增大的减速运动，则导体棒的加速度最大值出现在最低点或初始位置，导体棒速度最大时有 $mg=\frac{B^{2}L^{2}v}{R}+kx'$，

若没有磁场，由能量守恒定律有 $mgx=\frac{1}{2}kx^{2}$，

可知，在最低点时，有 $kx=2mg$，

此时，加速度为 $a=\frac{kx-mg}{m}=g$，

有磁场时，由能量守恒定律可知，弹簧的形变量减小，则最低点时加速度小于 $g$，则导体棒的加速度最大值出现在初始位置，此时，导体棒只受重力，加速度为 $g>\frac{B^{2}L^{2}v}{Rm}$，故 D 错误。

故选 C。
}

\item
%% number: 7
%% paperName: 2026年陕晋青宁高考第 7 题 4 分
%% typeId: 1
%% type: 单选题
%% chapter: 电场
%% point: 带电粒子的直线运动
%% method: 无
%% score: 4
%% degree: 450
%% duplicateId: 0
%% body:
飞行时间质谱仪可用于创新药物研发等领域，其简化工作原理如图。开关 S 闭合，在真空室内，电容器极板 AB 间加速电压为 $U$，两极板中央开有小孔，C 为接收端，BC 间为无场漂移区域。有一质量为 $m$、电荷量为 $q$（$q>0$）的带电大分子（不计重力），自漂入（忽略初速度）A 极板小孔起，经加速区域和漂移区域到达 C 端结束，时间探测器显示整个过程时间为 $t$，则 \xzanswer{C}

\onepicture[h=4.0cm]{\includesvg{figs/07}}

\fourchoices[answer=C]
{保持 S 闭合，仅 $\frac{q}{m}$ 变大，$t$ 变长}
{保持 S 闭合，仅 $U$ 增大，$t$ 变长}
{保持 S 闭合，AC 位置不变，仅 B 极板向右移动，$t$ 变长}
{将 S 断开，AC 位置不变，仅 B 极板向右移动，$t$ 变长}

%% answer: C
%% memo:
\memoanswer{
ABC．根据题意，设 A、B 距离为 $d$，B、C 距离为 $s$，且 $d+s=L$，在 A、B 间有 $\frac{qU}{d}=ma$，$qU=\frac{1}{2}mv^{2}$，$v=at_{1}$，

解得 $a=\frac{qU}{md}$，$v=\sqrt{\frac{2qU}{m}}$，$t_{1}=d\sqrt{\frac{2m}{qU}}$，

在 B、C 间有 $t_{2}=\frac{s}{v}=s\sqrt{\frac{m}{2qU}}$，

则整个过程时间

$t=t_{1}+t_{2}=d\sqrt{\frac{2m}{qU}}+s\sqrt{\frac{m}{2qU}}=(d+L)\sqrt{\frac{m}{2qU}}$。

保持 S 闭合，$U$ 不变，仅 $\sqrt{\frac{q}{m}}$ 变大，$t$ 变短，故 A 错误；

B．结合 A 分析可知，保持 S 闭合，仅 $U$ 增大，$t$ 变短，故 B 错误；

C．结合 A 分析可知，保持 S 闭合，$U$ 不变，AC 位置不变，仅 B 极板向右移动，则 $d$ 增大，$s$ 减小，$L$ 不变，$t$ 变长，故 C 正确；

D．将 S 断开，极板上电荷量 $Q$ 不变，AC 位置不变，仅 B 极板向右移动，A、B 间电场强度 $E=\frac{U}{d}=\frac{Q}{Cd}=\frac{4\pi kQ}{\varepsilon S}$ 不变，则加速度不变，则有 $v^{2}=2ad$，$t_{1}=\frac{d}{\frac{v}{2}}=\frac{2d}{v}$，$t_{2}=\frac{s}{v}$，

则整个过程时间

$t=t_{1}+t_{2}=\frac{2d+s}{v}=\frac{d+L}{\sqrt{2ad}}=\frac{1}{\sqrt{2a}}\left(\sqrt{d}+\frac{L}{\sqrt{d}}\right)$，

由于 $d$ 增大但小于 $L$，$s$ 减小，$L$ 不变，根据数学知识可知，当 $d=L$ 时，$\left(\sqrt{d}+\frac{L}{\sqrt{d}}\right)$ 有最小值，可知，随着 $d$ 增大但小于 $L$，$\left(\sqrt{d}+\frac{L}{\sqrt{d}}\right)$ 减小，则总时间 $t$ 减小，故 D 错误。

故选 C。
}

\item
%% number: 8
%% paperName: 2026年陕晋青宁高考第 8 题 6 分
%% typeId: 2
%% type: 多选题
%% chapter: 光学
%% point: 光的折射
%% method: 无
%% score: 6
%% degree: 550
%% duplicateId: 0
%% body:
玻璃镜片折射率是影响眼镜厚度的主要因素之一。如图，一束单色光从空气斜射入厚度为 $d$ 的平行玻璃砖左表面，从右表面射出，出射光线相对入射光线侧移量为 $\Delta x$。现换用折射率较大的玻璃砖，保持入射角不变，下列说法正确的是 \xzanswer{BC}

\onepicture[h=4.5cm]{\includesvg{figs/08}}

\fourchoices[answer=BC]
{若厚度 $d$ 相同，侧移量变小}
{若厚度 $d$ 相同，侧移量变大}
{若侧移量 $\Delta x$ 相同，厚度变小}
{若侧移量 $\Delta x$ 相同，厚度变大}

%% answer: BC
%% memo:
\memoanswer{
AB．根据题意，画出光路图（蓝色），现换用折射率较大的玻璃砖，保持入射角不变，折射角变小，若厚度 $d$ 相同，画出光路图（黄色），如图所示。

\onepicture[h=5.0cm]{\includesvg{figs/08_an1}}

由图可知，侧移量变大，故 A 错误，B 正确；

CD．根据题意，画出光路图（蓝色），现换用折射率较大的玻璃砖，保持入射角不变，折射角变小，若侧移量 $\Delta x$ 相同，画出光路图（黄色），如图所示。

\onepicture[h=5.0cm]{\includesvg{figs/08_an2}}

由图可知，厚度变小，故 C 正确，D 错误。

故选 BC。
}

\item
%% number: 9
%% paperName: 2026年陕晋青宁高考第 9 题 6 分
%% typeId: 2
%% type: 多选题
%% chapter: 机械波
%% point: 波的图象
%% method: 无
%% score: 6
%% degree: 450
%% duplicateId: 0
%% body:
一列简谐横波在均匀介质中沿直线传播，a、b 为介质中平衡位置相距 $0.5\Um$ 的两个质点。$t=0$ 时刻的波形图如图，$0.5\Us$ 时 a 振动至波谷，$1.2\Us$ 时 b 第一次回到波峰，则 \xzanswer{AC}

\onepicture[h=3.8cm]{\includesvg{figs/09}}

\fourchoices[answer=AC]
{波的周期为 $1.2\Us$}
{波速为 $10\Ums$}
{波的振幅为 $4\sqrt{3}\Ucm$}
{$0\sim0.6\Us$ 内质点 a 的路程为 $8\sqrt{2}\Ucm$}

%% answer: AC
%% memo:
\memoanswer{
A．根据题意由图可知，$t=0$ 时，质点 b 在波峰，$1.2\Us$ 时 b 第一次回到波峰，则波的周期为 $1.2\Us$，故 A 正确；

B．由于 $t=0.5\Us=\frac{5}{12}T<\frac{1}{2}T$ 时 a 振动至波谷，则 $t=0$ 时，质点 a 沿 y 轴负方向振动，则波沿 x 轴正方向传播，质点 a 由波峰振动到此位置的时间为 $\frac{1}{12}T$，则质点 a 的振动传播到质点 b 需要 $t'=\frac{1}{12}T=0.1\Us$，则波速为 $v=\frac{x}{t'}=\frac{0.5}{0.1}\Ums=5\Ums$，故 B 错误；

CD．设波的振幅为 $A$，结合上述分析可知，质点 a 的振动方程为 $y=A\sin\left(\frac{2\pi}{T}t+\frac{2}{3}\pi\right)\Ucm$，

又有 $t=0$ 时，$y=6\Ucm$，

解得 $A=4\sqrt{3}\Ucm$，

$\Delta t=0.6\Us=\frac{1}{2}T$，

经过 $t=0.6\Us=\frac{1}{2}T$ 时间，质点 a 通过的路程为 $s=2A=8\sqrt{3}\Ucm$，故 C 正确，D 错误。

故选 AC。
}

\item
%% number: 10
%% paperName: 2026年陕晋青宁高考第 10 题 6 分
%% typeId: 2
%% type: 多选题
%% chapter: 运动和力
%% point: 变加速直线运动
%% method: 无
%% score: 6
%% degree: 350
%% duplicateId: 0
%% body:
如图，物块 N、P、Q 的质量分别为 $m$、$m$、$2m$，Q 通过不可伸长的轻绳绕过两个固定轻质光滑定滑轮与 P 连接，P 与位于正下方的 N 用劲度系数为 $k$ 的轻弹簧相连。初始时托住 N 和 Q，使弹簧处于原长，三个物块均静止。现同时无初速度释放 N 和 Q，运动中，物块均视为质点且不与滑轮相碰，不计空气阻力，弹簧始终在弹性限度内，弹性势能 $E_{\mathrm{p}}=\frac{1}{2}kx^{2}$（$x$ 为形变量），重力加速度大小为 $g$。则 \xzanswer{AD}

\onepicture[h=4.5cm]{\includesvg{figs/10}}

\fourchoices[answer=AD]
{释放后瞬间 Q 的加速度大小为 $\frac{g}{3}$}
{释放后瞬间轻绳上的拉力大小为 $mg$}
{释放后 N 下降的最大距离为 $\frac{3mg}{4k}$}
{释放后 N 的速度最大值为 $\frac{g}{2}\sqrt{\frac{3m}{k}}$}

%% answer: AD
%% memo:
\memoanswer{
AB．根据题意可知，初始时弹簧处于原长，释放后瞬间，弹簧形变量不变，设释放后瞬间轻绳上的拉力大小为 $T$，对 Q、P 整体，由牛顿第二定律有 $2mg-mg=(m+2m)a$，对 Q，由牛顿第二定律有 $2mg-F=2ma$，

联立解得 $a=\frac{g}{3}$，$F=\frac{4}{3}mg$，故 A 正确，B 错误；

CD．根据题意，把物块 N、P、Q 看成整体，左边受向下的 $2mg$，右边也受向下的 $2mg$，整体合力为零，由动量守恒定律有 $mv_{N}=(m+2m)v_{PQ}$，

则有 $mv_{N}t=(m+2m)v_{PQ}t$，

可得 $mx_{N}=(m+2m)x_{PQ}$，

解得 $v_{N}=3v_{PQ}$，$x_{N}=3x_{PQ}$。

设释放后 N 下降的最大距离为 $x_{N}=h$，则 $x_{PQ}=\frac{1}{3}h$，此时弹簧形变量为 $\Delta x_{1}=h+\frac{1}{3}h$，

由能量守恒定律有 $mgh+2mg\cdot\frac{1}{3}h-mg\cdot\frac{1}{3}h=\frac{1}{2}k(\Delta x_{1})^{2}$，

解得 $h=\frac{3mg}{2k}$。

释放后 N 的速度最大时为 $v_{N}=v_{m}$，则有 $v_{PQ}=\frac{1}{3}v_{N}=\frac{1}{3}v_{m}$，

又有 $k\Delta x_{2}=mg$，

解得 $\Delta x_{2}=\frac{mg}{k}$，

则有 $x_{N}'=\frac{3mg}{4k}$，$x_{PQ}'=\frac{mg}{4k}$，

由能量守恒定律有

$mgx_{N}'+2mg\cdot x_{PQ}'-mg\cdot x_{PQ}'=\frac{1}{2}k(\Delta x_{2})^{2}+\frac{1}{2}mv_{m}^{2}+\frac{1}{2}(m+2m)\left(\frac{1}{3}v_{m}\right)^{2}$，

解得 $v_{m}=\frac{g}{2}\sqrt{\frac{3m}{k}}$，故 C 错误，D 正确。

故选 AD。
}

\item
%% number: 11
%% paperName: 2026年陕晋青宁高考第 11 题 5 分
%% typeId: 4
%% type: 实验题
%% chapter: 直线运动
%% point: 打点计时器公式
%% method: 无
%% score: 5
%% degree: 650
%% duplicateId: 0
%% body:
某同学使用打点计时器进行“测量速度和加速度”的实验。
\begin{enumerate}
	\item
	下述操作正确的是 \tkanswer{B}（填正确答案标号）。

	\twochoices[answer=B]
	{先拉动纸带，再开始打点}
	{先开始打点，再拉动纸带}
	\item
	该同学用电源频率为 $50\UHz$ 的打点计时器从左至右打出如图所示的一条纸带。可用某点左右相邻点间的平均速度代表该点的瞬时速度，则 P 点的瞬时速度为 \tkanswer{$0.305$} $\Ums$（保留三位有效数字）。
	\item
	测得 Q 点的瞬时速度 $v_{Q}=0.560\Ums$，则 P、Q 两点间的平均加速度为 \tkanswer{$1.28$} $\Umsq$（保留三位有效数字）。
\end{enumerate}
\onepicture[w=13.0cm, align=right]{\includesvg{figs/11}}

%% answer:
\jdanswer{
\begin{enumerate}
	\item
	B

	\item
	$0.305$

	\item
	$1.28$
\end{enumerate}
}
%% memo:
\memoanswer{
（1）使用打点计时器时，为了使打点稳定并有效利用纸带，应先接通电源开始打点，待打点稳定后再释放纸带拉动其运动。若先拉动纸带再打点，可能导致纸带上打出的点很少甚至没有点，且初始阶段打点不稳定。

故选 B。

（2）由图可知，P 点左侧相邻点对应的刻度为 $x_{1}=2.40\Ucm$，右侧相邻点对应的刻度为 $x_{2}=3.62\Ucm$。

P 点左右相邻两点间的距离 $\Delta x=x_{2}-x_{1}$，

打点计时器电源频率为 $50\UHz$，打点周期 $T=0.02\Us$，根据用某点左右相邻点间的平均速度代表该点的瞬时速度，可得 P 点的瞬时速度 $v_{P}=\frac{\Delta x}{2T}=0.305\Ums$。

（3）根据加速度定义式，P、Q 两点间的平均加速度

$a=\frac{v_{Q}-v_{P}}{10T}=\frac{0.560-0.305}{0.2}\Umsq\approx1.28\Umsq$。
}

\item
%% number: 12
%% paperName: 2026年陕晋青宁高考第 12 题 10 分
%% typeId: 4
%% type: 实验题
%% chapter: 电路
%% point: 电路实验
%% method: 无
%% score: 10
%% degree: 450
%% duplicateId: 0
%% body:
在电阻阻值测量中常常会用到平衡比较法。某实验小组测量待测电阻 $R_{x}$ 的阻值。
\begin{enumerate}
	\item
	实验步骤如下：

	①图甲中，开关 $S_{1}$ 断开时灵敏电流表 G（零刻线在表盘正中间）的指针指向正中间，开关 $S_{1}$ 闭合时 G 指针向 A 端偏转。用图乙所示的电路进行实验，开关 $S_{2}$ 闭合前，滑动变阻器 $R$ 的滑片应该置于 \tkanswer{D} 端（填“C”或“D”）。

	②图乙中，a、b、c 是三个完全相同的电阻，$R_{M}$ 为电阻箱，$R_{x}$ 为待测电阻。闭合开关 $S_{2}$，调节滑动变阻器 $R$，发现灵敏电流表 G 指针向 B 端偏转，则此时应调节电阻箱使 $R_{M}$ 阻值 \tkanswer{增大}（填“增大”或“减小”），使 G 指零（指针指向正中间）。

	③读得电阻箱 $R_{M}$ 的阻值为 $846\UO$，则待测电阻 $R_{x}$ 的阻值为 \tkanswer{$423$} $\UO$。

	\item
	当 a、b、c 三个电阻不完全相同时：

	①用图乙所示电路按上述步骤使灵敏电流表 G 指零，读得 $R_{M}$ 的阻值记为 $R_{1}$，在不增加实验器材的前提下需增加实验步骤完成实验，该步骤为 \tkanswer{A}（单选，填正确答案标号）。

	\threechoices[answer=A]
	{交换“a”和“b、c”的位置，调节电阻箱使 G 重新指零，读 $R_{M}$ 阻值}
	{将滑动变阻器的滑片向 C 滑动，使 G 重新指零，读 $R_{M}$ 阻值}
	{交换“$R_{M}$”和“b、c”的位置，调节电阻箱使 G 重新指零，读 $R_{M}$ 阻值}

	②按所选步骤读得 $R_{M}$ 的阻值记为 $R_{2}$，则待测电阻 $R_{x}$ 的阻值为 \tkanswer{$\sqrt{R_{1}R_{2}}$}（用 $R_{1}$、$R_{2}$ 表示）。
\end{enumerate}
\twopicture[numstyle=chinese, labelprefix={}, h=4.5cm, align=right]
{figs/12a.pdf}{figs/12b.pdf}

%% answer:
\jdanswer{
\begin{enumerate}
	\item
	①D；②增大；③$423$

	\item
	①A；②$\sqrt{R_{1}R_{2}}$
\end{enumerate}
}
%% memo:
\memoanswer{
（1）①为了保护灵敏电流表 G，闭合开关 $S_{2}$ 前，应使测量电路（电桥）两端的电压最小，故滑片应置于 D 端。

②由题意可知，电流从 A 端流入时指针向 A 端偏转。图乙中 G 指针向 B 端偏转，说明 B 端电势高。要使 G 指零，需增大 $\varphi_{A}$ 或减小 $\varphi_{B}$。根据电路图可知，增大 $R_{M}$ 可使 $R_{x}$ 分压减小，$\varphi_{A}$ 升高，故应调节电阻箱使 $R_{M}$ 阻值增大。

③图乙为电桥电路，平衡时满足 $\frac{R_{x}}{R_{M}}=\frac{R_{a}}{R_{bc}}$，

已知 a、b、c 为三个完全相同的电阻，代入平衡条件解得 $R_{x}=423\UO$。

（2）①当 a、b、c 不完全相同时，为了消除电阻阻值不准确带来的系统误差，可采用交换法。即交换比例臂电阻的位置。原平衡条件为 $\frac{R_{x}}{R_{1}}=\frac{R_{a}}{R_{bc}}$，

交换 a 和 b、c 的位置后，新平衡条件为 $\frac{R_{x}}{R_{2}}=\frac{R_{bc}}{R_{a}}$，

两式相乘可消去 $R_{a}$ 和 $R_{bc}$，解得 $R_{x}=\sqrt{R_{1}R_{2}}$，

故选 A。

②由上述分析可知 $R_{x}=\sqrt{R_{1}R_{2}}$。
}

\item
%% number: 13
%% paperName: 2026年陕晋青宁高考第 13 题 10 分
%% typeId: 6
%% type: 计算题
%% chapter: 气体性质
%% point: 玻意耳定律
%% method: 无
%% score: 10
%% degree: 650
%% duplicateId: 0
%% body:
某小组设计并完成了“稳定平抛水柱”实验。如图，竖直放置的储水瓶液面上方封闭少量气体，底部竖直细管甲与大气连通、水平细管乙有一阀门 K。初始时 K 关闭，瓶内上方气体体积 $V_{1}=200\UmL$，甲管内恰无水且 A 端口距水面高 $h_{1}=20\Ucm$，乙管 B 端口距水面高 $h_{2}=30\Ucm$；打开 K，水持续从乙管流出，大气通过甲管进入瓶内，当水面与甲管 A 端口齐平时关闭 K，此时瓶内上方气体的总体积 $V_{2}=1510\UmL$。已知大气压强 $p_{0}=1.00\times10^{5}\UPa$、水的密度 $\rho=1.00\times10^{3}\Ukgmc$，重力加速度 $g$ 取 $10\Umsq$，忽略温度变化与瓶体形变，气体均可视为理想气体。求：
\begin{enumerate}
	\item
	初始 K 关闭时，瓶内上方气体压强 $p_{1}$ 和乙管 B 端口处压强 $p_{2}$；
	\item
	在打开 K 到关闭 K 的过程中，进入瓶内的空气在大气压强下的体积 $\Delta V$。
\end{enumerate}
\onepicture[h=6.0cm, align=right]{\includesvg{figs/13}}

%% answer:
\jdanswer{
\begin{enumerate}
	\item
	$9.8\times10^{4}\UPa$，$1.01\times10^{5}\UPa$

	\item
	$1314\UmL$
\end{enumerate}
}
%% memo:
\memoanswer{
（1）初始 K 关闭时，甲管与大气连通且管内恰无水，说明 A 端口处压强等于大气压强 $p_{0}$。

由液体压强公式可知，瓶内上方气体压强 $p_{1}$ 满足 $p_{1}+\rho gh_{1}=p_{0}$，

代入数据解得 $p_{1}=9.8\times10^{4}\UPa$。

乙管 B 端口处压强 $p_{2}$ 满足 $p_{2}=p_{1}+\rho gh_{2}$，

代入数据解得 $p_{2}=1.01\times10^{5}\UPa$。

（2）当水面与甲管 A 端口齐平时，瓶内上方气体压强等于大气压强 $p_{0}$。设进入瓶内的空气在大气压强下的体积为 $\Delta V$，对瓶内原有气体和进入的空气整体应用玻意耳定律，有 $p_{1}V_{1}+p_{0}\Delta V=p_{0}V_{2}$，

代入数据解得 $\Delta V=1314\UmL$。
}

\item
%% number: 14
%% paperName: 2026年陕晋青宁高考第 14 题 13 分
%% typeId: 6
%% type: 计算题
%% chapter: 运动和力
%% point: 牛顿定律的应用
%% method: 无
%% score: 13
%% degree: 450
%% duplicateId: 0
%% body:
如图，某游乐场有一条滑道，由两段粗糙程度不同的直道组成，其中 $AB$ 段的长度 $l=36\Um$、倾角为 $\theta$，$BC$ 段水平且足够长。初始时游客甲乘滑板从 $A$ 点由静止下滑，经过 $t=6\Us$ 到达 $B$ 点，此后进入 $BC$ 段继续滑行。游客甲和滑板的总质量 $m=60\Ukg$，$AB$ 段、$BC$ 段滑道与滑板间的动摩擦因数分别为 $\mu_{1}$、$\mu_{2}$，其中 $\mu_{2}=\frac{1}{10}$。重力加速度 $g$ 取 $10\Umsq$，$\sin\theta=\frac{1}{3}$。游客和滑板整体视为质点，不计空气阻力及在 $B$ 点处的机械能损失。
\begin{enumerate}
	\item
	求 $AB$ 段滑道与滑板间的动摩擦因数 $\mu_{1}$；
	\item
	求 $AB$ 段甲和滑板的机械能损失及动量变化量的大小；
	\item
	当甲到达 $B$ 点时，游客乙乘同样滑板恰以 $5\Ums$ 的速度经过 $A$ 点下滑。求乙到达 $B$ 点时与甲的距离。
\end{enumerate}
\onepicture[w=11.0cm, align=right]{\includesvg{figs/14}}

%% answer:
\jdanswer{
\begin{enumerate}
	\item
	$\mu_{1}=\frac{\sqrt{2}}{10}$

	\item
	$2880\UJ$，$720\Ukgms$

	\item
	$40\Um$
\end{enumerate}
}
%% memo:
\memoanswer{
（1）当甲从 A 点滑至 B 点根据 $L=\frac{1}{2}at^{2}$，

代入数据可得 $a=2\Umsq$。

根据牛顿第二定律有 $mg\sin\theta-\mu_{1}mg\cos\theta=ma$，

又因为 $\sin\theta=\frac{1}{3}$，所以 $\cos\theta=\frac{2\sqrt{2}}{3}$，

代入数据联立可得 $\mu_{1}=\frac{\sqrt{2}}{10}$。

（2）甲从 A 点滑至 B 点损失的机械能等于这个过程中摩擦力做的功，所以有

$\Delta E=W_{f}=\mu mg\cos\theta\cdot L=2880\UJ$，

根据前面小问可知甲下滑的加速度为 $a=2\Umsq$；根据运动学公式 $v^{2}=2aL$ 可得甲到达 B 点的速度为 $12\Ums$，

所以 AB 段动量的变化量为 $\Delta p=mv-0=720\Ukgms$。

（3）乙从 A 滑到 B 的过程根据运动学公式有 $L=v_{0}t'+\frac{1}{2}at'^{2}$，

代入数据解得 $t'=4\Us$。

甲在地面上做匀减速运动，根据牛顿第二定律有 $\mu_{2}mg=ma_{2}$，

可得 $a_{2}=1\Umsq$。

甲到达 B 点的速度为 $v=12\Ums$，在 $4\Us$ 内甲没有停止；

在 $t'=4\Us$ 内，甲运动的距离 $x=vt'-\frac{1}{2}a_{2}t'^{2}$，

代入数据解得 $x=40\Um$，即甲乙之间的距离为 $40\Um$。
}

\item
%% number: 15
%% paperName: 2026年陕晋青宁高考第 15 题 16 分
%% typeId: 6
%% type: 计算题
%% chapter: 磁场
%% point: 带电粒子在组合场中的运动
%% method: 无
%% score: 16
%% degree: 350
%% duplicateId: 0
%% body:
如图，垂直于 $x$ 轴的足够大绝缘薄挡板 P 紧贴平行金属板放置，挡板 P 有一小孔在原点 $O$ 处，$O$ 点到垂直于 $y$ 轴的两金属板距离相等。在 $xOy$ 平面位于 $(-L,0)$ 处的粒子源以速率 $v_{0}$ 向面内各方向发射质量为 $m$、电荷量为 $q$（$q>0$）的带电粒子。两金属板（上板为正极）间距为 $2L$，板间电压 $U$ 初始为 $\frac{48mv_{0}^{2}}{25q}$。在 $xOy$ 平面 $y$ 轴右侧，圆心 $O'$ 位于 $(0,L)$、直径在 $y$ 轴、半径为 $L$ 的半圆区域存在方向垂直于平面向里的匀强磁场，其磁感应强度大小为 $\frac{mv_{0}}{qL}$。在 $x>L$ 处有一垂直于 $x$ 轴的固定弹性绝缘挡板 Q，粒子若碰到 Q 将发生弹性碰撞且电荷量 $q$ 不变。粒子打到金属板或挡板 P 均会被吸收，不计粒子重力及粒子间的相互作用，忽略边缘效应。（$v_{0}$、$m$、$q$、$L$ 均为已知量）
\begin{enumerate}
	\item
	可从小孔进入磁场的粒子，求其从 $O$ 点处射入磁场速度方向与 $x$ 轴正向夹角的正弦值；
	\item
	粒子从磁场射出后与挡板 Q 发生碰撞，求其再次射入磁场到达挡板 P 上位置的 $y$ 坐标；
	\item
	若仅改变电压 $U$ 的大小（$U>0$），在不同电压下，粒子源发射的所有粒子均无法打到挡板 P 右侧面；进入磁场的所有粒子打到挡板 P 右侧面有 1 个撞击点或 2 个撞击点，求以上三种情况分别对应的电压取值范围（用已知量 $v_{0}$、$m$、$q$ 表示）。
\end{enumerate}
\onepicture[h=7.5cm, align=right]{\includesvg{figs/15}}

%% answer:
\jdanswer{
\begin{enumerate}
	\item
	$0.6$ 或 $0.8$

	\item
	$\frac{2}{5}L$ 或 $\frac{4}{5}L$

	\item
	$U>\frac{2mv_{0}^{2}}{q}$ 时，无粒子能到达 $O$ 点，均无法打到 P 右侧面；当 $0<U<\frac{16mv_{0}^{2}}{17q}$ 或 $U=\frac{2mv_{0}^{2}}{q}$ 时，1 个撞击点；当 $\frac{16mv_{0}^{2}}{17q}\leq U<\frac{2mv_{0}^{2}}{q}$ 时，2 个撞击点。
\end{enumerate}
}
%% memo:
\memoanswer{
（1）粒子在电场中做类斜抛运动

\onepicture[w=7.5cm]{figs/15_an1.jpg}

设粒子初速度方向与 $x$ 轴夹角为 $\theta$，加速度

$a=\frac{qU}{m(2L)}=\frac{24v_{0}^{2}}{25L}$，

方向沿 $y$ 轴正向。粒子从 $(-L,0)$ 运动到 $(0,0)$，水平方向 $L=v_{0}\cos\theta\cdot t$，

竖直方向 $0=-v_{0}\sin\theta\cdot t+\frac{1}{2}at^{2}$，

联立解得 $v_{0}^{2}\sin(2\theta)=aL=\frac{24}{25}v_{0}^{2}$，

$\sin(2\theta)=\frac{qU}{2mv_{0}^{2}}=\frac{24}{25}$。

在 $O$ 点，$v_{x}=v_{0}\cos\theta$，$v_{y}=-v_{0}\sin\theta+at=v_{0}\sin\theta$，

设射入磁场速度方向与 $x$ 轴夹角为 $\beta$，则 $\tan\beta=\frac{v_{y}}{v_{x}}=\tan\theta$，

即 $\beta=\theta$。

根据 $\sin(2\beta)=\frac{24}{25}$，$\sin^{2}\beta+\cos^{2}\beta=1$，

解得 $\sin\beta=0.6$ 或 $\sin\beta=0.8$。

（2）粒子在磁场中做匀速圆周运动，半径 $R=\frac{mv_{0}}{qB}=L$，

磁场区域半径也为 $L$。由几何关系，粒子从 $O$ 点射入，必从磁场边界另一点射出，且射出速度方向水平向右，如下图

\onepicture[w=7.5cm]{figs/15_an2.jpg}

射出点纵坐标 $y_{\text{出}}=L(1-\cos\alpha)$，

当 $\sin\alpha=\frac{3}{5}$ 时，$\cos\alpha=\frac{4}{5}$，$y_{\text{出}}=\frac{L}{5}$；

当 $\sin\alpha=\frac{4}{5}$ 时，$\cos\alpha=\frac{3}{5}$，$y_{\text{出}}=\frac{2L}{5}$。

粒子与 Q 板弹性碰撞后沿 $x$ 轴负方向返回，再次进入磁场。

\onepicture[w=7.5cm]{figs/15_an3.jpg}

由对称性及圆周运动几何关系，粒子再次打到挡板 P 上的位置纵坐标 $y=2y_{\text{出}}$，

故 $y$ 坐标为 $\frac{2}{5}L$ 或 $\frac{4}{5}L$。

（3）根据（1）可知 $k=\frac{qU}{2mv_{0}^{2}}$。粒子能到达 $O$ 点需满足 $\sin(2\theta)=k$，即 $k\leq1$，$U\leq\frac{2mv_{0}^{2}}{q}$，

若 $U>\frac{2mv_{0}^{2}}{q}$，则 $k>1$，无粒子能到达 $O$ 点，均无法打到 P 右侧面。

当 $U\leq\frac{2mv_{0}^{2}}{q}$ 时，存在两个发射角 $\theta_{1}$、$\theta_{2}$（假设 $\theta_{2}$ 大于 $\theta_{1}$）。粒子在电场中不撞板条件为 $|y_{\min}|\leq L$，如下图

\onepicture[h=3.5cm]{figs/15_an4.jpg}

右图为刚好到达上极板的临界，根据逆向看，类平抛运动速度偏转角是位移偏转角正切值的 2 倍，可知，此时 $\tan\theta=2\times\frac{L}{\frac{1}{2}L}=4$，

根据数学关系可知此时 $\sin\theta=\frac{4\sqrt{17}}{17}$，$\cos\theta=\frac{\sqrt{17}}{17}$，

解得对于 $\theta_{2}$ 需 $k=\sin(2\theta)\geq\frac{8}{17}$，即 $U\geq\frac{16mv_{0}^{2}}{17q}$，

此时 $\theta_{2}$ 对应粒子轨迹刚好与上极板相切。

综上所述，当 $0<U<\frac{16mv_{0}^{2}}{17q}$ 时，仅 $\theta_{1}$ 粒子能进入磁场，1 个撞击点；

当 $\frac{16mv_{0}^{2}}{17q}\leq U<\frac{2mv_{0}^{2}}{q}$ 时，两粒子均能进入且撞击点不同，2 个撞击点；

当 $U=\frac{2mv_{0}^{2}}{q}$ 时，两粒子轨迹重合，1 个撞击点；

当 $U>\frac{2mv_{0}^{2}}{q}$，则 $k>1$，无粒子能到达 $O$ 点，均无法打到 P 右侧面。
}

\end{enumerate}

\end{document}
